% Transcription — fonds Alexandre Grothendieck, shelfmark 29, batch 2 (pages 21-40).
% Established from the facsimile archives/batches/29/batch-02.pdf (a mirror of
% https://grothendieck.umontpellier.fr/29.pdf), per the skill
% transcribe-grothendieck.
%
% Every one of the twenty pages carries mathematics; none is skipped.
%
% Pages 21-33 finish the typescript begun at page 12 of batch 1, "Données de
% ramification sur un schéma". Page 20 broke off at "Soit donnée une section
% du"; page 21 resumes at "faisceau associé au préfaisceau précédent". The
% batch carries the end of §1 (the kummerian datum attached to a divisor, the
% two normal-crossings cases, and the doctrine that a ramification datum is
% auxiliary — what is studied is the cartesian functor r), then §2 on the
% exactness properties of Rev^R(S) and the group pi_1^R(S,s), §3 on
% functoriality in S and in R, §4 on Galois coverings, and finally the
% "Remarques concernant l'exemple 3 du par. 1" of pages 32-33, which replace
% that example's artificial data by five conditions a) to e) on R alone.
%
% Page 31 changes language in mid-page: the French sentence stating the
% necessary and sufficient condition ends, and Grothendieck goes on in English
% to say why the category, and not the covering, is the right object. Both stay
% as written.
%
% Pages 34-38 are the covering letter, in English, that goes with the
% typescript — an answer to Murre's complaint that he had not grasped the
% set-up. It ends "I will probably drop by on June 2", which is the day Murre
% announced for Bures in his letter of 16 May 1967 (page 3 of batch 1). The
% letter itself bears no date and none is supplied here.
%
% Page 39 is the sheet Grothendieck headed "Finitude (Pbs et conjectures)" in
% ink; page 40 is the first sheet under it, a hard, heavily struck manuscript
% listing three finiteness statements for pi_1, pi_1' and pi_1^t over an
% algebraically closed field, with the state of proof of each. The scan is
% 228 ppi, which is the ceiling on how much of it can be read; several words
% under the strikings are left as \ill{}.
%
% The typescript's arrows and Greek letters were never inked on this carbon:
% the machine did not carry them and the gaps were meant to be filled by hand.
% They are restored here where the text settles them, and marked where it does
% not — one note at the first occurrence, as in batch 1.
%
% The dating is the inventory's own, for the whole shelfmark.
%
% Pass: Opus 5 (claude-opus-5), 2026-08-22 — first pass, unchecked against the pages by a human.
\documentclass[11pt,a4paper]{article}
% Le préambule commun à toutes les transcriptions du fonds Grothendieck.
%
% Il tient en une idée : **l'incertitude se compose, elle ne se résout pas.**
% Une transcription qui lisse un mot douteux a détruit la seule chose pour
% laquelle on la faisait — un lecteur doit pouvoir savoir, à chaque endroit,
% ce qui était sur le feuillet et ce qui a été ajouté. Les six macros qui
% suivent sont donc l'appareil critique, et non de la décoration.
%
% Les mêmes commandes sont reconnues par `scripts/render.mjs`, qui produit la
% vue de lecture du site. Une macro ajoutée ici doit l'être aussi là-bas,
% faute de quoi le rendu échouera bruyamment — ce qui est voulu : un rendu
% silencieusement faux est pire qu'un rendu absent.

\ProvidesPackage{grothendieck}[2026/08/08 Transcriptions du fonds Grothendieck]

\usepackage{amsmath,amssymb,amsthm}
\usepackage{mathrsfs}
\usepackage{stmaryrd}
% Les diagrammes commutatifs sont partout dans ce fonds : c'est la langue
% même de la géométrie algébrique de Grothendieck, et une transcription qui
% les rendrait en prose ne transcrirait rien.
\usepackage{tikz-cd}
% Français par défaut — c'est la langue des feuillets — mais l'anglais est
% chargé aussi : la traduction et le résumé sont des documents anglais, et la
% typographie française (espaces avant les deux-points) y serait une faute.
% Ces documents basculent par \selectlanguage{english} en tête de corps.
\usepackage[english,french]{babel}
\usepackage{fontspec}
\usepackage{geometry}
\geometry{a4paper,margin=25mm,marginparwidth=28mm}
\usepackage{marginnote}
\usepackage{xcolor}
% ulem, not soul. `soul`'s \st and \ul are built on a character-by-character
% scan that XeLaTeX's font machinery breaks: the document fails with an
% undefined control sequence rather than degrading. ulem does the same two jobs
% and is Unicode-engine safe. `normalem` stops it hijacking \emph, which the
% transcriptions use for Grothendieck's own emphasis.
\usepackage[normalem]{ulem}
\usepackage{hyperref}
\hypersetup{colorlinks=true,linkcolor=black,urlcolor=black,pdfborder={0 0 0}}

% --- Métadonnées du lot -----------------------------------------------------
% Reprises telles quelles par le script de rendu, qui les affiche en tête de
% la vue de lecture. Elles fixent ce que le fichier prétend couvrir : sans
% elles, une transcription flotte sans point d'attache dans un fonds de seize
% mille feuillets.

\newcommand{\folder}[1]{\gdef\grothFolder{#1}}
\newcommand{\batch}[1]{\gdef\grothBatch{#1}}
\newcommand{\leaves}[2]{\gdef\grothFirst{#1}\gdef\grothLast{#2}}
\newcommand{\foldertitle}[1]{\gdef\grothFoldertitle{#1}}
\gdef\grothFolder{?}\gdef\grothBatch{?}\gdef\grothFirst{?}\gdef\grothLast{?}\gdef\grothFoldertitle{}

% --- Le résumé --------------------------------------------------------------
%
% Il ouvre la lecture modernisée, sous le titre « Résumé », et il est composé
% en petit. Ce n'est pas un ornement : c'est le seul passage du document qui
% ne soit pas des mathématiques, et un lecteur doit pouvoir le reconnaître
% avant de l'avoir lu — puis le sauter s'il connaît déjà le sujet, ou s'y
% arrêter s'il ne le connaît pas. Le filet qui le suit marque la reprise du
% corps.

\newenvironment{resume}
  {\small\setlength{\parskip}{0.45em}}
  {\par\medskip\hrule\medskip}

% --- Mots-clés ---------------------------------------------------------------
%
% En anglais, en clôture du résumé de la lecture modernisée : le vocabulaire
% moderne sous lequel chercher ce que les feuillets construisent. Le manifeste
% du site (`npm run manifest`) les extrait pour étiqueter la cote entière —
% cette ligne est la source unique des tags, il n'y a pas de fichier à côté.
\newcommand{\keywords}[1]{\par\smallskip\noindent
  {\footnotesize\textsc{Keywords} --- \textit{#1}}\par}

% --- L'appareil critique ----------------------------------------------------

% Le numéro de feuillet, en marge. C'est l'ancre qui permet de régler un
% désaccord sur une formule en regardant *un* feuillet plutôt qu'une chemise
% de six cents. Le site s'en sert aussi pour tourner les pages du fac-similé
% quand on fait défiler la transcription.
\newcommand{\leaf}[1]{%
  \marginnote{\footnotesize\color{gray!70}\textsf{#1}}%
  \label{leaf:#1}\ignorespaces}

% Illisible : on ne devine pas. Un mot inventé qui se lit comme les autres est
% le pire résultat possible.
\newcommand{\ill}{\textcolor{red!60!black}{[\,\ldots\,]}}

% Lecture douteuse : le mot est donné, et signalé comme incertain.
\newcommand{\uncertain}[1]{\uline{#1}}

% Ajout de l'éditeur — une lettre manquante, un mot sous-entendu.
\newcommand{\add}[1]{\textcolor{blue!50!black}{[#1]}}

% Biffé par Grothendieck lui-même. Ce qu'il a rayé fait partie du document :
% une rature dit souvent où la pensée a changé de direction.
\newcommand{\struck}[1]{\sout{#1}}

% Note du transcripteur, sur l'état matériel du feuillet.
\newcommand{\note}[1]{\par\smallskip{\small\color{gray!60!black}\textit{[#1]}}\par\smallskip}

% Note marginale de Grothendieck, distincte du corps du texte.
\newcommand{\marginal}[1]{\par\smallskip\noindent\hspace{1em}%
  {\small\color{gray!60!black}#1}\par\smallskip}

% --- Titre ------------------------------------------------------------------

\AtBeginDocument{%
  \begin{center}
    {\large\bfseries Fonds Alexandre Grothendieck --- cote n\textsuperscript{o}~\grothFolder}\\[2pt]
    {\itshape\grothFoldertitle}\\[4pt]
    {\small feuillets \grothFirst--\grothLast\ (lot \grothBatch)}\\[2pt]
    {\footnotesize\color{gray!60!black}Transcription établie d'après le fac-similé numérisé
     par l'Université de Montpellier.\\
     Les lectures incertaines sont soulignées, les illisibles notées [\,\ldots\,],
     les ajouts entre crochets.}
  \end{center}
  \vspace{1em}}

\endinput


\folder{29}
\batch{2}
\pages{21}{40}
\dating{1959-1969}
\watermark{Édition de démonstration}
\foldertitle{Groupe fondamental [Autour de SGA 4 et SGA 7] : lettres (1959, 1967, 1969), tapuscrits et copies de tapuscrit annotés (s.d.), notes manuscrites (s.d.)}

\begin{document}

\section*{« Données de ramification sur un schéma », suite (pages 21 à 33)}

\page{21}
faisceau associé au préfaisceau précédent (NB la condition canularesque mise
sur les $(D_i)$ assure que l'on a bien un préfaisceau). Alors localement pour
la topologie étale, on a donc une situation $(D_i)$ comme dans a). Cela suffit
pour définir une donnée de ramification $\underline{R}$ sur $B$, généralisant
d'ailleurs celle envisagée dans a) (du moins lorsque la famille
$(D_i)_{i \in I}$ de a) satisfait à la condition supplémentaire locale qu'on
vient de formuler). Je vois deux cas particuliers ayant de l'intérêt,
\struck{\ill{}} et \struck{\ill{})} où $\underline{C}$ peut être défini en
termes d'un diviseur $D$ donné sur $S$ :
\note{« sur $B$ » sur la page, pour « sur $S$ » ; et le $\underline{C}$ de la
dernière ligne est laissé en blanc, la machine ne portant pas les lettres
soulignées de ce tapuscrit --- comme les flèches et les lettres grecques, il
est restitué ici d'après le texte}

1\textsuperscript{o}) $S$ est localement noethérien, $D$ est à croisements
normaux, on prend comme $D_i$ les composantes irréductibles locales régulières
de $D$ (qui ne sont définies que localement pour la topologie étale).

2\textsuperscript{o}) $S$ est au dessus d'un préschéma $T$, $D$ est à
croisements normaux relativement à $T$, on prend comme $D_i$ les composantes
locales lisses sur \uncertain{$T$} (qui ne sont définies que localement pour la
topologie étale).
\note{la lettre est surfrappée sur la page ; le relativement à $T$ de la ligne
précédente demande $T$}

Dans le cas 1\textsuperscript{o}), la catégorie
$\mathrm{Rev}^{\underline{R}}(S)$ est équivalente à une sous-catégorie pleine
de la catégorie des revêtements de $S$, et sauf erreur (il faudrait le
vérifier) il en est de même dans le cas 2\textsuperscript{o}). Les revêtements
en question s'appellent, dans le\struck{ur} cas 1\textsuperscript{o}),
\emph{revêtements modérément ramifiés au sens large} relativement au diviseur
$D$ (NB on vérifie en effet que c'est une notion plus large que celle de
« revêtement modérément ramifié par rapport à $D$ », (signifiant : par rapport
à la famille réduite à $D$). Le théorème d'Abhyankar permet de caractériser ces
revêtements par la condition d'être étales sur \struck{\ill{}}
$S - \mathrm{supp}\,D$, normaux en les points au dessus de $\mathrm{supp}\,D$,
toute composante irréductible en domine une de $S$, enfin ils sont modérément
ramifiés au sens classique en les points maximaux de $D$. Dans le cas
2\textsuperscript{o}, il y aurait lieu d'appeler les revêtements obtenus
« modérément ramifiés pour $D$, relativement à la

\page{22}
base $T$ ». Lorsque $T$ est régulier, les notions 1\textsuperscript{o} et
2\textsuperscript{o} coïncident. Le théorème de spécialisation pour le groupe
fondamental tame peut se démontrer, par la technique de SGA 1, quand on utilise
la notion de revêtement modérément ramifié au sens large, relativement au
diviseur \struck{dans les $\underline{s}_T$} $D_\xi$ ($D$ étant à croisement
normaux relativement à $T$, $S$ propre et lisse sur $T$).

\emph{Remarques.} Evidemment, pour tout $S'$ objet de $\underline{Et}_S$, on
définit de façon évidente $\mathrm{Rev}^{\underline{R}}_{\ill{}}(S')$, en
considérant la restriction de $\underline{R}$ à $\underline{Et}_{S'} =
(\underline{Et}_S)_{/S'}$. Pour $S'$ variable, on obtient ainsi une catégorie
fibrée $\underline{\mathrm{Rev}}^{\underline{R}}$ sur $\underline{Et}_S$, dont
les fibres $\mathrm{Rev}^{\underline{R}}_{S'}$ sont les
$\mathrm{Rev}^{\underline{R}}(S')$, et un foncteur cartésien
\[
  (*) \qquad \underline{r} : \underline{\mathrm{Rev}}^{\underline{R}}
  \longrightarrow \underline{\mathrm{Rev}} ,
\]
fibrée sur $\underline{Et}_S$, où on désigne pour abréger par
$\underline{\mathrm{Rev}}$ la catégorie des revêtements des différents objets
$S'$ de $\underline{Et}_S$. (NB on suppose que pour tout $S'$, le foncteur
réalisation géométrique sur $\underline{C}_{S'}$ est à valeurs dans les
revêtements de $S'$, et que le foncteur réalisation géométrique sur
$\mathrm{Rev}^{\underline{R}}(S')$ est bien défini). Il faut insister sur le
fait que pour nous, la « donnée de ramification » $\underline{R}$ est un objet
de nature essentiellement auxiliaire, destiné à définir la catégorie fibrée
$\underline{\mathrm{Rev}}^{\underline{R}}$ (qui est même un \emph{champ}) et le
foncteur cartésien $\underline{r}$ de $(*)$ ; c'est $(*)$ qui est l'objet
véritable de l'étude. On regardera donc deux données de ramification
$\underline{R}$, $\underline{R}'$ comme essentiellement équivalentes, si elles
définissaient des foncteurs $(*)$ équivalents (par une équivalence, bien sûr,
induisant l'identité sur $\underline{\mathrm{Rev}}$). On peut se convaincre sur
des exemples que $\underline{R}$ et $\underline{R}'$ peuvent avoir des allures
très différentes. Ainsi, sous les conditions de l'exemple 3), si
$\underline{R}'$ est défini par une sous-catégorie fibrée $\underline{C}'$ de
$\underline{C}$

\page{23}
telle que pour tout objet $M$ d'un $\underline{C}_{S'}$, il existe un
recouvrement de $S'$ par des $S'_i$ tel que pour tout $i$, $M_{S'_i}$ soit
\struck{dominé} majoré par un $M_i$ objet de $\underline{C}'_{S'_i}$,
\struck{\ill{}} et supposant que le foncteur réalisation géométrique pour
$\underline{R}'$ est celui induit pour le foncteur réal.\ géom.\ pour
$\underline{R}$, --- alors on constate que $\underline{R}$ et $\underline{R}'$
sont « équivalents ». Comme cas particulier, signalons l'exemple de la donnée
kummérienne a), dans le cas où les $D_i$ admettent des équations globales par
des sections $c_i$ de $\underline{O}_S$. Alors on peut considérer le système
projectif des $Z^{\underline{n}}_{\underline{c}}$, pour $\underline{n} =
(n_i)_{i \in I}$ variables, premiers \struck{aux caractéristiques résiduelles}
à la caractéristique de $S$ (supposée exister). Considérons de plus le système
projectif des groupes $Z_{\underline{n}} = \mathbb{Z}/n_i\mathbb{Z}$, et
supposons donné un isomorphisme entre celui-ci et le système projectif des
$\mu_{\underline{n}}$ (il en existe lorsque $S$ est un schéma sur un corps
algébriquement clos, par exemple). On peut alors considérer les
$Z^{\underline{n}}_{\underline{c}}$ comme des schémas à groupes
$Z_{\underline{n}}$, formant un système projectif, \struck{\ill{}} donc on est
sous les conditions de l'exemple 2, définissant une donnée de ramification
$\underline{R}'$ sur $S$. On constate alors que cette donnée est équivalente à
la donnée $\underline{R}$ décrite dans le cas particulier kummérien a).

Comme autre exemple de données de ramification équivalentes, signalons le
suivant, où $\underline{R}$ est la donnée de ramification considérée dans
l'exemple 1. On se donne un revêtement principal $P$ de $S$, de groupe $G'$
(opérant à droite), et on suppose que $G$ opère dessus librement à gauche, en
commutant aux opérations de $G'$. Soit alors $Z' = Z \times^G P$, considéré
comme revêtement à groupe d'opérateurs $G'$, d'où une donnée de ramification
$\underline{R}'$. Cette dernière est équivalente à $\underline{R}$, comme on
voit en considérant les équivalences \struck{quasi-inverses \ill{}} de
\struck{\ill{}} \struck{entre} $\underline{\mathrm{Rev}}^{\underline{R}}$ dans
$\underline{\mathrm{Rev}}^{\underline{R}'}$ \ill{} opérateurs

\page{24}
évidente lorsque $P = G'_S$ avec l'opération régulière droite de $G'$, et
l'opération gauche de $G$ est donnée par un monomorphisme $G \to G'$ ; on se
ramène à ce cas particulier par localisation étale.

On peut se proposer de normaliser la notion de donnée de ramification, pour
éliminer l'arbitraire dans le choix d'un $\underline{R}$. Mais alors il semble
qu'on soit infailliblement conduit à prendre comme catégorie $\underline{C}$ la
catégorie $\mathrm{Rev}^{\underline{R}}$ elle-même qu'on se propose de
construire et d'étudier ! (Comparer avec la situation analogue pour les
relations entre sites et topos, l'objet auxiliaire étant le site, l'objet
intrinsèque de nature géométrique étant le topos associé à ce site).

\textbf{2. Propriétés de la catégorie $\mathrm{Rev}^{\underline{R}}(S)$. Le
groupe fondamental $\pi_1^{\underline{R}}(S,\ )$.}

On constate immédiatement que la catégorie $\mathrm{Rev}^{\underline{R}}(S)$
hérite de toutes les propriétés d'exactitude \struck{\ill{}} établies pour la
catégorie des revêtements étales d'un schéma donné, et dont une partie sont
énoncés dans Exp V n\textsuperscript{o}4 (conditions G1 à G3). De façon
précise, on peut dire que dans la catégorie $\mathrm{Rev}^{\underline{R}}(S)$,
les limites projectives finies et les limites inductives finies existent, et
satisfont aux mêmes propriétés d'exactitude que dans la catégorie des ensembles
(telle l'exactitude du foncteur changement de base dans cette catégorie) ; en
particulier tout morphisme de la catégorie se factorise en le produit d'un
épimorphisme et d'un monomorphisme, les épimorphismes et monomorphismes sont
\struck{\ill{}} effectifs ; enfin tout monomorphisme est un isomorphisme avec
un sommande direct. On fera attention que le foncteur « réalisation
géométrique » sur $\mathrm{Rev}^{\underline{R}}(S)$ commute bien aux limites
inductives finies (et \struck{\ill{}} en particulier aux sommes finies), mais
non pas en général aux limites projectives finies, et en particulier, pas aux
produits de deux objets.

Pour appliquer la théorie de Galois générale de V 4, il faut dis-

\page{25}
poser de plus d'un \emph{foncteur fibre}
\[
  F : \mathrm{Rev}^{\underline{R}}(S) \longrightarrow (\text{Ens finis}) ,
\]
(conditions G4, G5 de V 4) i.e.\ d'un foncteur exact. Dans le cas de l'exemple
1, tout point géométrique de $Z$ fournit un tel foncteur fibre, et ce foncteur
est conservatif (condition G6) si on suppose que toute partie ouverte et fermée
de $Z$ stable par $G$ est vide ou égale à $Z$. De même, dans le cas de
l'exemple 2, on obtient un \struck{\ill{}} foncteur fibre en prenant un point
géométrique de $Z = \varprojlim Z_i$. Pour obtenir un foncteur fibre dans le
cas plus général de l'exemple 3, on choisit un point géométrique $s$ de $S$
\struck{\ill{}} qui n'appartienne pas à \struck{\ill{}} « l'ensemble de
ramification » de $\underline{R}$, i.e.\ tel que pour tout $S'$ dans
$\underline{Et}_S$ et tout objet $M$ de $\underline{C}_{S'}$, le revêtement à
opérateurs (groupe $G$) $\underline{r}(M)$ \struck{\ill{}} de $S'$ soit
principal en les points de $S'$ au dessus de l'image de $s$ dans $S$. Alors
\struck{\ill{}} on constate que pour tout objet $T$ de
$\mathrm{Rev}^{\underline{R}}(S)$, $\underline{r}(T)$ est étale sur $S$ en $s$,
et le foncteur $T \mapsto \underline{r}(T)(s)$ (fibre géométrique en $s$) est
bien exact. Pour qu'il soit conservatif, il faut et il suffit que le site
$\underline{C}$ (pour la topologie image inverse de celle de $\underline{Et}_S$)
soit connexe ; \struck{\ill{}} si on suppose, avec les notations de l'exemple
3, que pour tout $M$, définissant $(\underline{r}(M), G)$,
$\underline{r}(M)/G \to S'$ soit un isomorphisme (ce qui sera sans doute le cas
dans tous les cas intéressants), alors il revient au même de dire que $S$ est
connexe. Appliquant la théorie de V 4, on trouve donc un groupe pro-fini, noté
$\pi_1^{\underline{R}}(S, s)$, \struck{\ill{}} et une équivalence de la
catégorie des ensembles finis où ce groupe opère continûment, avec la catégorie
$\mathrm{Rev}^{\underline{R}}(S)$.

\textbf{3. Fonctorialités pour $S$ et $\underline{R}$ variables.}

Considérons un morphisme $f : T \to S$ de schémas, d'où un foncteur

\page{26}
continu
\[
  f^{\circ} : \underline{Et}_S \longrightarrow \underline{Et}_T .
\]
On suppose donné une donnée de ramification $\underline{R}_S =
(\underline{C}_S, \underline{s}_S)$ sur $S$, une autre $\underline{R}_T =
(\underline{C}_T, \underline{s}_T)$ sur $T$, et un foncteur noté
$M \mapsto M \times_S T : \underline{C}_S \to \underline{C}_T$, transformant
morphismes cartésiens en morphismes carétsiens et rendant commutatif de
diagramme
\note{« carétsiens » et « rendant commutatif de diagramme » sur la page}

\begin{tikzcd}
\underline{C}_S \arrow[r] \arrow[d] & \underline{C}_T \arrow[d] \\
\underline{Et}_S \arrow[r, "f^{\circ}"] & \underline{Et}_T
\end{tikzcd}

Considérons de plus le carré de foncteurs

\begin{tikzcd}[column sep=huge]
\underline{C}_S \arrow[r] \arrow[d, "\underline{s}_S"'] & \underline{C}_T \arrow[d, "\underline{s}_T"] \\
(\mathrm{Sch})_{/S} \arrow[r, "X \mapsto X \times_S T"'] & (\mathrm{Sch})_{/T}
\end{tikzcd}

et supposons donné un morphisme $u$ entre les deux composés :
\[
  u(M) : \underline{s}_T(M \times_S T) \longrightarrow
  \underline{s}_S(M) \times_S T
\]
qui, pour tout $M$ fixé \struck{\ill{}} sur $S' \in \mathrm{Ob}\,\underline{Et}_S$,
soit un $f^{\circ}(S')$-morphisme. Sous ces conditions, on définit de façon
évidente un foncteur, noté $f^{\circ}$ si une confusion n'est pas à craindre
(ou $M \mapsto M \times_S T$), de $\mathrm{Rev}^{\underline{R}}(S)$ dans
$\mathrm{Rev}^{\underline{R}_T}(T)$ et un homomorphisme de foncteurs composés

\begin{tikzcd}[column sep=huge]
\mathrm{Rev}^{\underline{R}_{(S)}}(S) \arrow[r] \arrow[d] & \mathrm{Rev}^{\underline{R}_{(T)}}(T) \arrow[d] \\
\mathrm{Rev}(S) \arrow[r, "X \mapsto X \times_S T"'] & \mathrm{Rev}(T)
\end{tikzcd}

Cet homomorphisme est un isomorphisme dans certains cas, par exemple sous les
conditions de l'exemple 3, lorsqu'on suppose (sous les notations d'icelui) que
pour tout sous-groupe $H$ de $G$, le passage au quotient
$\underline{s}(M)/H$ commute au changement de base $T \to S$ ; cette condition
sera vérifiée

\page{27}
par exemple si $T \to S$ est plat, ou dans les cas kummériens a) envisagés plus
haut (en supposant que les images inverses des $D_i$ dans $T$ existent, et
prenant ces derniers pour définir une donnée de recollement sur $T$).

On recule devant la tâche d'expliciter une propriété de transitivité pour un
composé $U \to T \to S$, et jetant provisoirement un voile pudique sur ces
abîmes, nous allons remarquer plutôt que le foncteur
$\mathrm{Rev}^{\underline{R}_{(S)}}(S) \to
\mathrm{Rev}^{\underline{R}_{(T)}}(T)$ qu'on vient de construire a le bon goût
d'être \emph{exact}. Si alors, dans les conditions de l'exemple 3, $t$ est un
point \emph{géométrique} de $T$, non ramifié pour $\underline{R}_{(T)}$ dont
l'image dans $S$ est non ramifié pour $\underline{R}_{(S)}$, et si on suppose
que $u$ dans l'homomorphisme $(*)$ de foncteurs composés induit un isomorphisme
sur les fibres géométriques en $t$, alors les sorites de Exp V impliquent qu'on
a un homomorphisme canonique $\pi_1^{\underline{R}_{(T)}}(T, t) \to
\pi_1^{\underline{R}_{(S)}}(S, s)$. Ceci s'applique en particulier au cas de
données kummériennes type a), ou type b) 1\textsuperscript{o} et
2\textsuperscript{o}, quand on suppose que $T \to S$ est un morphisme régulier
resp.\ lisse (de sorte que les images inverses des $D_i$ satisfont encore aux
mêmes conditions de « normal crossings » que $D$ lui-même).

Il y a a expliciter deux cas particuliers importants de fonctorialité. L'un est
celui où $T \to S$ est étale, et où $\underline{R}_{(T)}$ est induit par
$\underline{R}_{(S)}$ sur $S$, auquel le \struck{foncteur} carré nous était
déjà connu par ailleurs (et l'homomorphisme de foncteurs composés y est un
isomorphisme), puisque l'on dispose d'un foncteur cartésien
$\underline{\mathrm{Rev}}^{\underline{R}} \to \underline{\mathrm{Rev}}$ sur
$\underline{Et}_S$. Lorsqu'on suppose de plus $T$ fini sur $S$, $S$ et $T$
connexes, et qu'on est dans les conditions de l'exemple 3, alors
l'homomorphisme sur les groupes fondamentaux est un monomorphisme, dont l'image
est le sous-groupe de $\pi_1^{\underline{R}}(S, s)$ défini par le revêtement
ponctué $T$ de $S$. Le deuxième cas est celui où $T \to S$ est un morphisme
entier, radiciel et surjectif, et où $\underline{R}_{(T)}$

\page{28}
est « l'image inverse » de $\underline{R}_{(S)}$, toujours sous les conditions
de l'exemple 3. Alors il est \emph{trivial} que
$\mathrm{Rev}^{\underline{R}_{(S)}}(S) \to
\mathrm{Rev}^{\underline{R}_{(T)}}(T)$ est une équivalence, utilisant le
résultat analogue pour les revêtements étales de schémas $X$ et
$X \times_S T$ (où $X$ \struck{est} $= \underline{s}(M)$ est un schéma
\struck{\ill{}} sur $S$ \dots). En particulier, quand on dispose d'un point
géométrique non ramifié de $T$, on obtient un isomorphisme sur les groupes
fondamentaux. Le même raisonnement, après localisation étale à peu près
évidente, s'applique au cas kummérien \struck{\ill{}}.

J'ai fait allusion plusieurs fois à la donnée de définir une donnée de
ramification image inverse d'une donnée de ramification sur $S$ par un
morphisme $T \to S$ ; la définition explicite serait nécessairement un peu
lourde ; je me borne à signaler que dans le cas où $\underline{R}$ est défini
par $(Z, G)$ comme dans l'exemple 1, son image inverse est définie par
$(Z_T = Z \times_S T, G)$. L'ennui avec cette définition générale, c'est
qu'elle ne s'applique pas directement au cas où $\underline{R}$ est une donnée
kummérienne associée à un système $(D_i)$, car la donnée image inverse au sens
précis précédent n'\emph{est pas} la donnée correspondante au système
$(D_{i\,T})$ même à équivalence de catégories fibrées près ; elle lui est
seulement « essentiellement équivalente » au sens du n\textsuperscript{o}1.
Seule une rédaction en forme permettra de se rendre compte s'il y a lieu
d'introduire cette notion d'image inverse, ou s'il ne vaut pas mieux donner des
conditions axiomatiques sur la donnée $(*)$, sous les conditions de l'exemple
3, qui assurent moralement que $\underline{R}_{(T)}$ est essentiellement
équivalente à l'image inverse de $\underline{R}_{(S)}$. Bien sûr, ces
conditions, pour \struck{\ill{}} être utiles, doivent s'appliquer directement
aux cas kummériens. Ce sera sans doute le parti le plus raisonnable. Il me
semble également qu'il serait raisonnable, assez rapidement après les
définitions général-

\page{29}
es, de se borner aux données de ramification $\underline{R}$ du type envisagé
dans l'exemple 3 ; ce sont les seules sans doute qu'on rencontrera jamais, et
les seules aussi où on sait construire le foncteur réalisation géométrique
$\underline{r}$ et en dire des choses raisonnables.

Avec l'une ou l'autre approche mentionnée à l'alinéa précédent pour traiter de
données de ramifications sur des schémas variables sur $S$, (et toujours sous
les conditions de l'exemple 3 \dots) on est en condition pour énoncer et
démontrer qu'un morphisme fidèlement plat et quasi-compact $T \to S$ est de
descente effective pour la catégorie fibrée sur
$(S, T, T \times_S T, T \times_S T \times_S T)$ formée des revêtements de type
$\underline{R}_{(S)}$ \dots\dots{} Ceci s'appliquera en particulier à la
descente fidèlement plate des revêtements kummériens ; c'est utilisé (sans
démonstration) à la page 38 des notes de Murre dans le cas de la descente à
partir d'un anneau complété, sans compter diverses utilisations de descente à
partir d'un localisé strict. Il y a également lieu d'énoncer une propriété de
passage à la limite, pour un système projectif $(X_i)$ de schémas
quasi-compacts et quasi-séparés à morphismes de transition affines ; cela
permettra de ramener ce qui se passe au voisinage d'un point $s$ d'un $S$ à
voir ce qui se passe en $\mathrm{Spec}(\underline{O}_{S,s})$. Par exemple, dans
le cas $S$ loc.\ noeth.\ normal muni d'une famille $(D_i)_{i \in I}$ de
diviseurs réduits, donc où la catégorie des revêtements kummériens s'identifie
à une sous-catégorie de la catégorie des revêtements, on trouve que le
revêtement $X$ de $S$ est \struck{\ill{}} kummérien \uncertain{ssi} pour tout
$s \in S$, il en est de même du revêtement induit sur
$\mathrm{Spec}(\underline{O}_{S,s})$, et on peut dans cet énoncé remplacer
$\underline{O}_{S,s}$ \struck{par} son hensélisé ou par son hensélisé strict,
\struck{\ill{}} voire par son complété si les hypothèses faites restent stables

\page{30}
par passage au complété (par exemple si $\underline{O}_{S,s}$ à fibres
formelles géom.\ normales, par exemple s'il est excellent). Les mêmes
observations s'appliqueront sans hypothèse de normalité ou de réduction
quelconque, à condition de se borner aux revêtements
\struck{\uncertain{kummériens}} mod.\ ram.\ \emph{galoisiens} ; cf
n\textsuperscript{o} suivant.

\textbf{4. Cas des revêtements galoisiens.}
\note{le chiffre du titre est surfrappé sur la page ; le n\textsuperscript{o}
précédent étant le 3, et le renvoi de l'alinéa qui précède visant « le
n\textsuperscript{o} suivant », c'est bien le 4}

Un revêtement \emph{galoisien} de type $\underline{R}$ de $S$, à groupe de
Galois $G$, est un objet $X$ de $\mathrm{Rev}^{\underline{R}}(S)$, sur lequel
le groupe fini $G$ opère à droite, de telle façon que dans la catégorie
$\mathrm{Rev}^{\underline{R}}(S)$ on obtienne un revêtement principal de
l'objet final, i.e.\ 1\textsuperscript{o}) le morphisme
$X \times G \to X \times X$ est un isomorphisme et 2\textsuperscript{o}) le
morphisme $X \to$ objet final est couvrant. Comme
$\mathrm{Rev}^{\underline{R}}(S)$ est en tous cas une catégorie
multigaloisienne, la notion de revêtement galoisien de type $\underline{R}$
s'explicite également en termes des groupes fondamentaux décrivant cette
catégorie. On peut enfin l'expliciter en revenant à la définition de
$\mathrm{Rev}^{\underline{R}}(S)$, et en exigeant que les revêtements des
$\underline{s}(M)$ ($M \in \mathrm{Ob}\,\underline{C}_{S'}$) correspondants à
$X$ soient principaux de groupe $G$. Par abus de langage, le revêtement à
opérateurs $(\underline{r}(X), G)$ de $S$ s'appellera également revêtement
galoisien de type $\underline{R}$ de $S$. Cet abus de langage ne présente pas
d'inconvénient trop sérieux dans le cas où le foncteur de la catégorie des
revêtements galoisiens de type $\underline{R}$ de $S$, dans la catégorie des
revêtements à opérateurs de $S$, défini par le foncteur « réalisation
géométrique », est pleinement fidèle, lorsqu'on se borne à prendre comme
morphismes dans la première les seuls isomorphismes. (Alors, en effet, pour un
revêtement à opérateurs $(T, G)$ donné de $S$, il existe, à isomorphisme
canonique près, au plus un $(X, G)$ \struck{\ill{}} galoisien de type
$\underline{R}$ dont il provient). Nous allons donner des conditions pour qu'il
en soit bien ainsi, en nous plaçant comme d'habitude dans le cas de l'exemple
3. Voici d'abord la condition néc.\ et suffisan-

\page{31}
te dans le cas particulier de l'exemple 1 : pour tout localisé strict $S'$ d'un
point de $S$, prenant une composante connexe $Z'_0$ de $Z' = Z \times_S S'$ et
son stabilisateur $H$ dans $G$, et pour deux sous-groupes distingués
$H_1, H_2$ de $H$ et tout isomorphisme $u : Z'_0/H_1 \to Z'_0/H_2$,
\struck{\ill{}} compatible avec un isomorphisme
$\varphi : H/H_1 \to H/H_2$, il existe un élément $g$ de $H$ tel que
$\mathrm{int}(g)$ applique $H_1$ sur $H_2$ et induise $\varphi$ par passage au
quotient, et tel que $u$ soit induit par l'endomorphisme $g_{Z'_0}$ de $Z'_0$
par passage aux quotients. Dans le cas général de l'exemple 3, il faut exiger
ceci pour tous les couples $(\underline{s}(M), G)$ sur les divers
$S' \in \underline{Et}_S$.

This condition is satisfied in the case of Kummer data, as one sees by
combining 1.10.\ and 1.15.\ of Murre's notes. This phenomenon explains why I
have been so reluctant (wrongly so) to define \struck{Kummer} tamely ramified
coverings except in the Galois case : in the non galois case, it will have to
be not just a covering satisfying some extra \emph{conditions}, but a covering
with some extra \emph{structure} ; whereas in the Galois case, the extra
structure is uniquely determined, if it exists. Note however that \struck{\ill{}}
not all morphisms (in the usual naive sense) between (even Galois) tame
coverings which are not isomorphisms should be considered, only those which
come form morphisms in $\mathrm{Rev}^{\underline{R}}(S)$, --- that one has to
be careful is shown by the counterexample to prop.1.16.\ of Murre's notes.
Also, the natural notion of products of tame \struck{\ill{}} coverings, whether
galois or not, is not the naïve fibered product over $S$ (which by the way
accounts for the complication of the proof of Murre's lemma 4.3.). This shows
rather clearly that even in the Galois case, the consideration of the category
$\mathrm{Rev}^{\underline{R}}(S)$ is undispensable for a good comprehension of
what is really going on.
\note{« come form », pour « come from », et « undispensable » sont sur la page}

\subsection*{Remarques concernant l'exemple 3 du par.\ 1 (pages 32 et 33)}

\page{32}
\emph{Remarques} concernant l'exemple 3 du par.1. --- La présentation de cet
exemple n'est ni élégante ni naturelle, la donnée de $\underline{s}$ et de
\ill{} étant un élément de structure artificiel. Il vaut mieux présenter les
choses ainsi, en termes de conditions sur $\underline{R}$, savoir les
suivantes :

\begin{itemize}
\item[a)] Pour tout objet $M$ d'un $\underline{C}_{S'}$, le faisceau étale
associé au préfaisceau sur $\underline{Et}_{S'}$ des automorphismes de $M$ est
\struck{\ill{}} représentable par un schéma en groupes $\underline{G}(M)$ fini
et étale sur $S$ (évidemment déterminé à isomorphisme unique près).
\end{itemize}

Il est évident que $G(M)$ opère sur le schéma relatif $\underline{s}(M) = Z$.

\begin{itemize}
\item[b)] Pour tout objet $S'$ de $\underline{Et}_S$, tout morphisme dans
$\underline{C}_{S'}$ est un épimorphisme, et tout endomorphisme d'un objet de
$\underline{C}_{S'}$ est un automorphisme. Pour tout morphisme
$u : M' \to M$ dans $\underline{C}_{S'}$, et tout automorphisme $g'$ de $M'$,
il existe un automorphisme $g$ de $M$ tel que $gu = ug'$.
\end{itemize}

Cet automorphisme est évidemment unique, de plus la condition précédente
implique qu'à $u : M' \to M$ est associé un unique homomorphisme de schémas en
groupes $G(M') \to G(M)$ tel que $u$ soit compatible avec les opérations de
$G(M')$ et $G(M)$ sur $M'$ et $M$.

\begin{itemize}
\item[c)] Avec les notations précédentes, \struck{si} l'homomorphisme
$G(M') \to G(M)$ est surjectif, et si $H$ désigne son noyau, le
\struck{\ill{}} morphisme $\underline{s}(M') \to \underline{s}(M)$ induit un
\struck{morphisme} isomorphisme $\underline{s}(M')/H \to \underline{s}(M)$,
\struck{en fait de \ill{}} \struck{un fibré principal sur $\underline{s}(M)$
de groupe} (NB on suppose que les $\underline{s}(M)$ sont \struck{relativement
des \ill{}} relativement affines, de sorte que le passage au quotient a un
sens).

\item[d)] Deux homomorphismes de $M'$ dans $M$ sont, localement pour la
topologie étale, déduits l'un de l'autre par composition avec un automorphisme
de $M$.
\end{itemize}

\page{33}
\begin{itemize}
\item[e)] Etant donnés deux objets $M, M'$ de $\underline{C}_{S'}$, il existe
une famille couvrante de morphismes $S'_i \to S'$ dans $\underline{Et}_S$,
telle que pour tout $i$, les images $M_i$ et $M'_i$ de $M$ et $M'$ sur $S'_i$
soient dominés par un même objet $M''_i$ de $\underline{C}_{S'_i}$.
\end{itemize}

Lorsqu'on adopte cette présentation, on peut, dans la description des données
de ramification kummériennes, se dispenser d'introduire aussi des isomorphismes
\struck{\ill{}} $u : G_{S'} \to \mu_{\underline{n}}\,S'$.

\section*{Lettre à J.-P. Murre (pages 34 à 38)}

\page{34}
Dear Murre,

I am very sorry I did not succeed to convey the intuitive idea behind the
general nonsense of my notes. It seems to me that the basic example in order to
understand the idea is example 1, where you can take $Z$ to be a standard
Kummer covering for definiteness, $Z = Z^n_a$, and $S'$ (normal, say)
\struck{normal}. Intuitively, when you look at coverings of $S$ whose
ramification type is not worse than the one of $Z$ over $S$, you mean that the
normalized inverse image $Z'$ of $S'$ over $Z$ is étale. \struck{\ill{}} From
the birational point of view, assuming $S'$ connected and therefore
corresponding to a field extension $K'$ of the field of functions $K$ of $S$,
this means simply that $K'$ is isomorphic to a subextension of an extension
$L'$ of the function field $L$ of $Z$, unramified with respect to the model
$Z$ ; when $S$ (and whence $Z$) is strictly local, hence $L'=L$, this means
simply that $K'$ is isomorphic to a subextension of $L$, a very reasonable way
indeed of interpreting the intuitive fact that $S'/S$ has no worse ramification
than $Z/S$. (NB We assume of course now, for simplicity, that $Z$ is connected,
and moreover that $Z$ is the normalisation of $S$ in some finite Galois
extension $L$ of $K$). Note that anyhow, we want that the notion in question
should be of local nature for the étale topology, thus we can reduce always to
the strictly local case. Now by transport de structure the group $G$ operating
on $Z$ must operate on $S'$ (in a way compatible with the operations on $Z$),
and it is immediate that we recover $S'$ from the $G$-covering $Z'$ of $Z$ as
being simply $S'=Z'/G$. Thus the category of \struck{\ill{}} étale coverings of
$S$ we are interested in is just equivalent to the catégory of coverings of
$(Z,G)$. This latter catégory on the other hand is of a very standard type ;
the presence of $G$ does not prevent it to share all the properties of
catégories of étale coverings

\page{35}
as developped in Exp VI, and under a mild restriction in the nature of a
connectedness assumption, we can say that the category in question is
\struck{classified} completely described by a certain pro-finite group, the
so-called fundamental group of $(Z,G)$, which in fact is an extension of $G$ by
the usual fundamental group of $Z$. On the other hand, for these latter
considerations, all normality assumptions are completely superfluous. Their
main use was to insure that the functor $Z' \mapsto Z'/G$ from étale coverings
of $(Z,G)$ to coverings of $S$ is fully faithful ; this latter fact, on the
other hand, reduces at once to the case when $Z$ is a principal covering of
$S$, with group $G$ (because this condition will be satisfied when restricting
to a non empty open set $U$ of $S$, and on the other hand the restriction
functors \struck{\ill{}} to $U$ resp.\ $Z/U$ are fully faithful, because of
normality) ; but in this case descent theory tells us that we get in fact an
\emph{equivalence} of the category of \emph{étale} coverings of $S$, and the
catégory of étale coverings of $(Z,G)$ :

In most applications however, instead of having just one model for ramification
(one $Z^n_a$ say), we have an infinity of such (all $Z^n_a$ say, fixed $a$ and
variable $n$), and the condition on coverings $S'$ of $S$ we are interested in
is that the covering should (at least locally) have ramification no worse than
one of the $Z_i$. This situation is described in example two. On the other
hand, sometimes the ramification models are not defined globally (for instance
in the Kummer case, when we give \struck{\ill{}} divisors which are not
principal), only locally, in a way that among the set of models given over
various open sets (or more generally, objects of the étale site) there are
certain compatibility conditions, roughly that the restriction of any model
over $U$ to $U \cap V$ has ramification no worse than the restriction of some
model given over $V$. Such an example is

\page{36}
dealt with in example 3, which has the disadvantage of being somewhat involved
and artificial in its technical details, but has the merit of allowing explicit
construction of the functor « geometric realization ». Otherwise, the set up
outlined before the examples seems more natural, and giving a closer grasp of
the \struck{\ill{}} really essential features of what « ramification data »
should really mean. (By the way, one shoul rather say « category of
ramification data »).

You are of course right that in the definition of a Galois covering when I
assume that the coverings of the various $\underline{s}(M)$ are Galois étale
coverings with Galois group $G$, that beforhand an action of $G$ on $X$ was
given, which means that $G$ acts on each of the coverings $X(M)$ of
$\underline{s}(M)$ on a way compatible with the base changes
$\underline{s}(M') \to \underline{s}(M)$ stemming from arrows
$M' \to M$. In the various examples 1 to 3, this means that the action of $G$
commutes to étale localization on $S$ (i.e.\ base changes $S' \to S$ in the
étale site of $S$) \emph{and} to the action of the groups of automorphisms
given on the various $Z$'s.

Now let's come back to the case $S$ strictly local ; then (for coverings
defined over the whole of $S$) the general situation of example 3 reduces to
the one of example 2, and replacing the $Z_i$'s by suitable connected
components, we may assume them connected if we want. For a given $Z_i$, the
\struck{\ill{}} étale coverings of $(Z_i, G_i)$ are trivial as coverings of
$Z_i$ alone (i.e.\ forgetting the action of $G_i$) because $Z_i$ is strictly
local too. Therefore the catégory of étale coverings of $(Z_i, G_i)$ is just
equivalent to the category of finite sets $E$ on which $G_i$ opérates ; to such
an $E$ is associated the constatnt covering $E_{Z_i}$ of $Z_i$, but $G_i$
operating on it via it's operation on $E$ ; the geometric relaization over $S$
of this is nothing else but $Z_i \times^G E$, the covering associated to the
$G$-covering
\note{« constatnt », « relaization », « opérates » sont sur la page}

\page{37}
$Z_i$ of $S$ and the action of $G_i$ on $E$. When $S$ is again arbitrary, this
shows in a rather concrete way, locally (for étale topology), what it means,
for a covering $S'$ of $S$, to be associated to some $M$ : it means that
locally it is describable by a representation of an inertial group on a finite
set, in the obvious way ; and the extra structure on one $S'$ involved by
deducing it from some $S$, can be described by saying that we give such local
representations of $S'$, in such a way that they match together i.e.\ satisfy
some rather evident compatibility condition (This description makes a sense
whenever the functor « geometric realization » is faithful --- not necessarily
fully faithful). The analogous description holds for Galois coverings with
Galois group $G$, namely in the strictly local case, for a fixed $Z_i$, the
\struck{\ill{}} category of these is equivalent to the category of finite sets
$E$ on which $G_i$ operates on the left, $H$ on the right, $G_i$ and $H$
commuting, in such a way that $E$ is a principal homogeneous space (or, as we
say, a torsor) under $H$ ; up to isomorphism, they are classified by
$\mathrm{Hom}(G_i, H)/H$, group homomorphisms mod.\ interior automorphism.
Using this description, one easily proves the criterion I indicated for the
« geometric realization-functor » to be fully faithful on the category of
Galois coverings of type $\underline{R}$ ; indeed, as usual the question is a
local one for the étale topology, which allows one to reduce to the stricly
local case.

Your interpretation of the Kummer case in the final formulation of example 3 is
indeed the one I had in mind. Also, when I wrote $n' \geqslant n$, I meant of
course the order relation of divisibility (it may be convenient to introduce
this order relation explicitly, for simplicity of notations).

\page{38}
I realize that all the indications I have given you so far are extremely
sketshy, and as a consequence that I am charging you with a considerable amount
of work to put some sense and order into all that. Thus it is I, not you, who
should apologize for causing a lot of trouble ! I look forward with great
pleasure meeting you in Bures. As I am having some russian and chinese lessons
on fridays, I will probably drop by on June 2.

With best regards

\section*{« Finitude (Pbs et conjectures) » (pages 39 et 40)}

\page{39}
\emph{Finitude} (Pbs et conjectures)

\page{40}
\emph{Théorèmes de finitude pour le groupe fondamental}

1) Sur un corps algébriquement clos :

\begin{itemize}
\item[a)] $X$ propre$/k \implies \pi_1(X)$ de type fini

\item[b)] $X$ t.f.$/k \implies \pi'_1(X)$ de \emph{présentation finie}
\marginal{$'$ indique la partie première à la caractéristique}

\item[c)] $X$ t.f.$/k \implies \pi^t_1(X)$ de type fini
\marginal{$t$ indique que \uncertain{la} ramification \ill{} donne des
« \uncertain{diviseurs} » $H$ \struck{\ill{}} \struck{d'un corps} au corps de
fonctions du \ill{}}
\end{itemize}

[NB.\ c) plus général que a)]
\note{$\pi'_1(X)$ et $\pi^t_1(X)$ sont entourés d'un même trait, et une flèche
monte de $\pi^t_1(X)$ vers $\pi'_1(X)$}

\marginal{$X_i \subset X$ irréd.\ \uncertain{composant}}

a) démontré, [par \uncertain{descente} réduction au \uncertain{normal}, par
Bertini au cas d'une courbe, puis théorie de spéc.\ du gpe fond.\ et méthode de
\uncertain{transcendance} \ill{}]

\struck{\ill{}}

\struck{b) réduction au cas $X$ normal [\uncertain{démontré}] \ill{}}

b) le cas « type \emph{fini} » est prouvé

[réduction au cas normal par descente, puis au cas où singularités « jolies »
au sens Artin, on utilise alors la suite exacte des $\pi_1$ i.e.\ th.\ de
spécialisation du gpe fond., puis \ill{} \struck{\ill{}} au cas d'une courbe
normale, \ill{} au cas d'une courbe \ill{} où on ne connaît \ill{} par le th.\
de spéc.\ \ill{} \uncertain{transcendant}]

\uncertain{Quid} : ceci, on est ramené dans b) au cas $X$ \emph{normal}
[démontré] et \struck{\ill{}} [ni \uncertain{résolution}] au cas $X$
\struck{\ill{}} \emph{lisse}$/k$. \emph{Si} $X$ \struck{propre}
\struck{[\ill{}]} \ill{} \uncertain{résolution} au cas d'une surface. \ill{}
d'une surface \uncertain{projective} \ill{} \struck{\ill{}} \ill{}
\struck{complété} diviseurs (à normal crossings).
\note{le bas de la page est écrit jusqu'au bord de la feuille et la dernière
ligne y est coupée ; les mots portés sous les ratures ne sont pas lisibles à la
définition de la numérisation}

\end{document}
